\documentclass[11pt]{article}
\usepackage[margin=0.9in]{geometry}
\usepackage{booktabs}
\begin{document}
\section*{Common 2+2 returned-ordinate congruence}
\texttt{projective-law-common22-residual/v1}. Lean proves the cleared
identity $C^5N + SC^2X + T = U[-C^4X + C^3(-Cq_1-A) + C^2Q]$ over a
commutative ring. Source-derived integer polynomials meet five obligations
connecting the Rust return tuple to the theorem's ordinate definitions.
Changing the sign of the $s$ term on either side is rejected.

\begin{center}
\begin{tabular}{lrr}
\toprule
Obligation & Monomials & Equal\\
\midrule
r\_has\_C\_squared\_factor & 1369 & yes\\
qdelta\_matches\_source & 81 & yes\\
rdelta\_matches\_source & 1384 & yes\\
s\_matches\_lean & 1577 & yes\\
t\_matches\_lean & 11350 & yes\\
\bottomrule
\end{tabular}
\end{center}
\noindent The term-level CSV retains every expanded integer coefficient. This
cleared polynomial statement does not establish weighted normalization, the
residual support's curve equation, the principal-divisor step, exceptional
branches, or a cost change.
\end{document}
